\bookchapter{第四十七章：速查手冊與資源}

\chapterintro{本章乃一份可獨立使用之速查手冊。當汝記不清某個原則或須快速決策時，直接翻到此裡。建議汝將它當成「口袋卡片」：五大原則速查、工具決策樹、常見問題、學習資源，一頁不落。}

\booksection{一、思維模型速查卡}

\begin{center}
\begin{tabularx}{\linewidth}{|X|X|X|}
\hline
\textbf{原則} & \textbf{一句話} & \textbf{操作要點} \tabularnewline
\hline
Think Before Coding & 先想清楚再開口 & 回答三問：做何？結果長何樣？有何約束？ \tabularnewline
Simplicity First & 用最少代碼解決問題 & 每次只做一個功能，砍掉不必要之複雜度 \tabularnewline
Surgical Changes & 精準修改不越界 & 明確指定修改範圍，拒絕AI之「順手改一下」 \tabularnewline
Goal-Driven Execution & 說「做何」不說「何以做」 & 予目標不予步驟，使AI自主選擇路徑 \tabularnewline
Iterate \& Verify & 生成—驗證—修正循環 & 三層驗證：功能、邊界、架構 \tabularnewline
\hline
\end{tabularx}
\end{center}

\begin{center}\small\color{gray}五大思維原則速查卡\end{center}

\booksection{二、工具選擇決策樹}

當汝須選擇AI編程工具時，按照以下決策路徑選擇。注意：此不為「哪個最好」之排名，乃「哪個最適合汝」之判斷框架。

汝要做何？

零基礎，想快速體驗 選雲端 Builder (v0/Bolt/Replit)

日常開發，須代碼補全 選 IDE 內置 AI (Cursor/Trae)

大型項目，須批量操作 選終端 Agent (Claude Code)

快速原型驗證想法 選雲端 Builder 或 Trae SOLO

\begin{center}\small\color{gray}圖6：AI編程工具選擇決策樹\end{center}

\booksection{三、常見問題FAQ}

Q1：吾完全不會編程，能用AI編程嗎？

能。2026年之AI編程工具已經足夠強大，汝完全可用自然語言描述需求，使AI生成代碼。但有兩點建議：第一，從雲端Builder工具（如v0、Bolt）開始，它等不須配置開發環境；第二，邊用邊學基礎概念——不須精通語法，但汝須能「看懂」AI所寫代碼之大意，如此纔能有效驗收。

Q2：吾當學哪個編程語言？

取決於汝要做何。做網站前端——JavaScript/TypeScript；做數據處理或AI——Python；做企業級後端——Java或Go。但更重要之為：於AI編程時代，「語言」之重要性於下降。AI能助汝處理語法細節，汝須理解者乃「程序之結構」——何為函數、何為循環、何為數據結構。此諸概念乃跨語言者，一旦理解，換何語言皆快。

Q3：AI所寫之代碼有版權問題嗎？

截至2026年8月，此仍然為一個有爭議之法律問題。一般來說，AI生成之代碼於大多數司法管轄區不被視為有明確之版權所有者。但若汝用於商業項目，建議：(1) 瞭解汝所用AI工具之服務條款中關於知識產權之條款；(2) 對關鍵代碼進行人工審查與修改，使其具有汝之創造性投入；(3) 避免使AI直接複製可能有版權之開源代碼。

Q4：AI會取代程序員嗎？

短期內不會。AI所取代者不為「程序員」此職業，乃「逐行編碼」此工作。就像計算器無取代數學家，乃使數學家從繁瑣之手算中解放出來一樣。AI編程工具使程序員從「寫代碼」轉向「想問題」——此實際上提高了程序員之價值上限。未來之競爭力不於於「能寫多少行代碼」，而於於「能解決多複雜之問題」。

Q5：如何判斷AI之產出為否正確？

三層驗證法。第一層：功能驗證——它能不能跑？輸入輸出對不對？此乃最基本者。第二層：邊界驗證——空輸入、超長輸入、異常輸入會怎樣？此能發現隱藏之問題。第三層：架構驗證——代碼結構為否合理？有無硬編碼？為否便於擴展？此須汝有一定之技術判斷力。若汝無法判斷第三層，至少做到前兩層——能跑、邊界正常，就已比大多數初學者強了。

\booksection{四、推薦學習資源}

以下為截至2026年8月之推薦學習資源。請注意，互聯網資源連結可能變化，若某個連結失效，搜索其標題通常能尋到最新版本。

\begin{center}
\begin{tabularx}{\linewidth}{|X|X|X|}
\hline
\textbf{類型} & \textbf{資源} & \textbf{推薦理由} \tabularnewline
\hline
原始文獻 & Andrej Karpathy之CLAUDE.md & 五大原則之源頭，70行但字字珠璣 \tabularnewline
原始文獻 & Anthropic官方文檔（Claude Code） & 終端Agent之最佳實踐指南 \tabularnewline
教程 & Vibe Coding社區整理之Skills合集 & 社區實踐者總結之實戰技巧集合 \tabularnewline
教程 & Cursor官方文檔與YouTube頻道 & IDE內置AI之權威使用指南 \tabularnewline
社區 & GitHub上之awesome-claude-code & Claude Code之最佳實踐與插件集 \tabularnewline
社區 & Reddit之r/cursor社區 & 全球Cursor用戶之經驗分享與討論 \tabularnewline
基礎 & MDN Web Docs (developer.mozilla.org) & 前端開發最權威之參考文檔 \tabularnewline
基礎 & Python官方教程（docs.python.org） & Python學習之官方起點 \tabularnewline
\hline
\end{tabularx}
\end{center}

\begin{center}\small\color{gray}推薦學習資源清單\end{center}

\booksection{五、結語：工具於手，方法論於心}

回到本書之核心信念：工具每三個月迭代一次，方法論每三年才進化一次。於2026年8月之今日，AI編程工具之能力已經強大到令人驚歎——汝能用一句話創建一個網站，用一段描述生成一個全棧應用，用一條指令完成一個複雜重構。但工具越強大，方法論之重要性反而越突出——蓋工具之強大放大了汝之每一個決策。

若汝之需求清晰，AI能助汝高效實現；若汝之需求含混，AI也會高效地做出錯誤之東西。若汝遵循簡單至上，AI能助汝快速驗證想法；若汝貪多求全，AI也會助汝快速堆積技術債。若汝精準修改，AI能助汝外科手術般地優化代碼；若汝指令含混，AI也會助汝大刀闊斧地引入bug。

此乃本書故花大部分篇幅講方法論而非工具操作之故。不為因為工具不重要，乃因為方法論乃工具之「乘數」——方法論越好，工具之價值越大；方法論為零，工具再強也只為零乘以任何數。

汝不復獨自編程。汝乃統帥，AI 乃汝之軍團。汝思之精粗，定軍團之戰力。自今日始，以方法論武其思，以工具張其能——此 AI 編程之全部祕密也。

\booksection{參考文獻}

\begin{enumerate}
  \item Andrej Karpathy. (2025). Vibe Coding. Social media post introducing the concept of vibe-driven programming with AI.
  \item Andrej Karpathy. (2025). CLAUDE.md - Project guidelines for AI-assisted development. GitHub repository, 70+K stars.
  \item Anthropic. (2026). Claude Code: Terminal-native AI programming. Official documentation and best practices.
  \item Cursor Team. (2026). Cursor 3: The next generation of AI-first IDE. Release notes and documentation.
  \item OpenAI. (2026). Codex CLI: Bringing AI to the terminal. Technical documentation.
  \item Model Context Protocol (MCP) Specification. (2025-2026). Open protocol for AI-tool integration.
  \item Vibe Coding Community. (2026). Community-curated skills and best practices for AI-assisted programming.
  \item Spec-Driven Development (SDD) methodology. (2026). Emerging practices in specification-first AI programming.
\end{enumerate}
